\documentclass[11pt]{article}
\usepackage[a4paper,margin=26mm]{geometry}
\usepackage{fontspec}
\setmainfont{lmroman10-regular.otf}[BoldFont=lmroman10-bold.otf,ItalicFont=lmroman10-italic.otf,BoldItalicFont=lmroman10-bolditalic.otf]
\usepackage{amsmath,amssymb,microtype}
\usepackage[unicode,hidelinks]{hyperref}
\setlength{\parindent}{0pt}\setlength{\parskip}{7pt}
\hypersetup{pdftitle={Paper IV — Apparatus — Prompt 02},pdfauthor={Editorial apparatus}}
\begin{document}
{\LARGE Paper IV}\par
{\large Apparatus to printed pp. 65–86}\par
Prompts 01–02. French source edition and independent English translation.

\section*{Witness and editorial treatment}
The controlling witness is the attached exact-scope PDF, leaves 3–24, corresponding
to full-volume PDF pages 82–103. The original R27 and R28 bytes remain preserved
as unverified prior work; only source-replayed derivative text is accepted.
Pages 81–86 are transcribed directly from the authority. Pages 63–64 are copy matter, not author prose.

The French preserves the printed wording, notation, quoted statements, display
grouping, and page seams. The English preserves all arguments and mathematical
relations, including the legible source problems described below. No conjectural
repair is inserted into either text. Running heads and the imprint are translated
in the English reader. The author’s note remains a footnote; editorial comments
appear only here.

\section*{P. 66: unqualified converse}
In the paragraph beginning \emph{Il est facile de voir}, the print states a converse
from solvability for every prime factor of $B$ to solvability for $B$ itself.
The unrestricted statement is retained. It needs qualification when repeated prime
factors occur: for example, $A=3$, $B=9$ gives a solution modulo the prime factor $3$,
but no solution modulo $9$, whose fourth-power residues are $0,1,4,7$.
This counterexample is an editorial calculation, not text supplied by the witness.

Authority: exact-scope leaf 4 / full-PDF page 83, second paragraph.

\section*{P. 70: the printed plus sign}
In the paragraph beginning \emph{Le produit}, the print calls $K$ a divisor of
$\varphi'^2+2u'^2$. The plus sign is legible. Earlier on the same page, the factorization
uses $\varphi'^2-2u'^2$. The plus is reproduced in both editions. The surrounding
argument indicates a likely printed error; that diagnosis is editorial and does
not change the diplomatic reading.

Authority: exact-scope leaf 8 / full-PDF page 87, the paragraph before the quoted theorem.

\newpage
\section*{P. 71: restoration of the author’s note}
The inherited formula $r=\rho^2\mp1$ is an algebraic rewriting, not the printed
expression. The witness reads $r=b\sigma^2\pm1$, followed by $s=\rho\sigma$.
The two signs reduced to $+$ in R27’s norm identity are also printed as $\pm$:
\[
(rt\pm bsu)^2-b(ru\pm st)^2=p.
\]
Both editions restore the printed forms. The source marker is $*)$, attached after
\emph{impair} (\emph{odd} in English), before the sentence’s full stop. R27’s extra
closing parenthesis in the French body is not present in the source. The book title
\emph{Théorie des Nombres} is italic in the source; it is rendered as
\emph{Theory of Numbers} in English. The citation’s printed numbering is retained.

Authority: exact-scope leaf 9 / full-PDF page 88, complete footnote.

Detail witness: \path{qa/source/p071_footnote_detail.png}.

\section*{Other page records}
Pages 65, 67–69, and 72–74 require no conjectural emendation. The separate pagewise
apparatus files record the title and rules, the lexical reading \emph{de théorèmes}
on p.68, lower-case \emph{d’où} on p.69, the equation labels and repeated factor
arrays, and the exact continuation points. The cumulative catalogue has one record for each
of the twenty-two accepted author pages.

The sentence at the foot of p.67 resumes at the head of p.68; that at the foot of
p.71 resumes with the result labelled $(\alpha)$ on p.72. The reciprocity display
announced at the end of p.72 begins p.73. The substitution sentence at the foot of
p.73 resumes with its power of $u$ on p.74. No bridging prose is supplied.

Section 4 begins on p.74. Prompt 01 ended after its two-by-two factor display;
Prompt 02 resumes with the proof that $E$ and $F$ are relatively prime on p.75.
All earlier page sources remain byte-identical. The cumulative readers now stop
after the reciprocity row at the foot of p.86, not at the end of Paper IV.
Pages 87–98 are not accepted by this checkpoint; the final cold audit remains pending.

\section*{Layout record}
The readers are re-typeset, not photographic facsimiles. Original page boundaries,
paragraph starts, display membership, formula punctuation, and the author-note
anchor are retained. Physical within-page line wrapping and justification are not
copied; line-end word division is reflowed, without changing lexical hyphens such
as \emph{tout-à-fait} and \emph{à-peu-près}. These layout choices are explicit, not
silent lexical normalization. The untouched authority and page images remain the
reference for the original lineation and letterforms.

\newpage
\section*{P. 77: the printed variable $t$}
The paragraph beginning \emph{En combinant ce résultat} says that the outcome
depends on whether $t$ has the form $8n+7$ or $8n+3$. The letter is visibly $t$,
not $b$. Both editions retain it. The preceding equality is
$\left(\frac{t}{b}\right)=\left(\frac{2}{b}\right)$, while Theorem I on p.78 gives the
corresponding alternatives in terms of $b$. This discrepancy is recorded only
here; no substitution is made in the author text.

Authority: exact-scope leaves 15–16 / full-PDF pages 94–95.

\section*{P. 82: the printed square on $p$}
In the first paragraph, following \emph{On voit, par l'équation}, the source
reads $p^2=\varphi^2+\psi^2$. The superscript 2 is legible and is reproduced in
both editions. On p.81 the decomposition is introduced as $p=\varphi^2+\psi^2$.
The extra square is likely a printing error, but that diagnosis is editorial,
not a warrant for changing the diplomatic reading.

Authority: exact-scope leaf 20 / full-PDF page 99, first paragraph.

\section*{P. 85: $ps$, not $ps^2$}
The parenthetical replacement equation reads $t^2+au^2=ps$. There is no
superscript on this $s$, whereas the two immediately preceding equations
both have $ps^2$ on the right. The printed difference is retained in both
editions without supplying a conjectural exponent.

Authority: exact-scope leaf 23 / full-PDF page 102, the long methodological paragraph.

\section*{New page records and seams}
No conjectural emendation is required on pp.75–76, 78–81, 83–84, or 86. The
source equation labels $(\delta),(\delta'),(\varepsilon),(\zeta),(\eta),(\theta)$ and
$(\alpha')$ retain their local scope; the renewed $(\alpha')$ on p.86 belongs to the
\emph{Addition}. No renumbering is imposed.

The phrase ending p.77 continues with \emph{de $b$} on p.78. The factor list at
the foot of p.79 ends with $h,$ and resumes with $h',h'',\ldots$ on p.80.
The two sign cases on p.80 are joined in $(\eta)$ on p.81. The sentence at the
foot of p.83 continues with \emph{même chose} on p.84. The \emph{Addition au mémoire
précédent} heading and both adjacent rules are on p.85; its opening sentence
continues with \emph{simplifiées} on p.86. No bridging prose is inserted.

There are no new author footnotes on pp.75–86. The note mentioned on p.86
refers to the note near the beginning of section 3 on p.71, preserved in the
cumulative readers. Source signatures $10^*$, $11$, and $11^*$ occur on
pp.75, 81, and 83 respectively; p.81 also carries the volume imprint.

\end{document}
